\section{侦查终结、侦查羁押期限、补充侦查}
\begin{enumerate}
  \item 侦查终结。
        \begin{enumerate}
          \item 事实清楚；证据确实、充分；法律手续完备。
          \item 辩护人提出要求的应当听他的意见。被动听。
          \item 后续处理。
                \begin{enumerate}
                  \item 移送同级检察院审查起诉。
                  \item 撤销案件。救济：只能提起自诉（可以附带民事诉讼）。
                \end{enumerate}
        \end{enumerate}
  \item 侦查羁押期限。最长可以逮捕多长时间，\zd{2 + 1 + 2 + 2 + X}。
        \begin{enumerate}
          \item 一般 2 个月。
          \item 案情复杂，上一级检察院，+1 月。
          \item 交集流广重大复杂案件，省一级检察院，+2 月。
          \item 10 年以上，省一级检察院，+2 月。凡是延长 2 月的都得省一级检察院批准。
          \item 特殊原因，最高检报全国人大常委，+X 月。
          \item 不包括逮捕之前的先行拘留的时间。
          \item 公安机关重新计算侦查羁押期限的，应当报批准逮捕的检察院备案。
        \end{enumerate}
  \item \zd{补充侦查}。4 种常见的理由。\zd{事实不清，证据不足，漏罪漏人}。都是检做决定。\zd{不计入审理期限}。
        \begin{enumerate}
          \item 审查批捕阶段。\yi 说明理由。\er 通知公安补侦（\scx{退回}）。
          \item 审查起诉阶段。
                \begin{enumerate}
                  \item \yi \zd{捕诉部门}自己来 \hz 退回。哪里来哪里去。自侦 \hz 监委。退回优先于自己来。
                  \item 每次 1 月，\zd{2 次为限}，不管管辖转、变。
                  \item 退 1/2 次补侦 \yjt 仍然证据不足 \yjt 可以/应当不起诉。
                \end{enumerate}
          \item 法庭审理/审判阶段。\zd{除了审查起诉阶段，其余都不能退回补侦}。
                \begin{enumerate}
                  \item 被告人提供了新的立功线索，法院\zd{可以建议}检察院启动补侦。\zd{只有建议权，无退回权}。
                  \item 补侦主体。原则检察院（已经指控的事实）\hz 要求公安（漏罪）补侦。\zd{不是退回}。
                  \item 期限、次数同审查起诉阶段。
                \end{enumerate}
        \end{enumerate}
\end{enumerate}
